\chapter{شيفرة مورس}
% full version: chapters/04_code.tex

في شيفرة مورس رمزان، النقطة والشرطة، وفواصل بينهما. أما أبجدية العصبون فأفقر: رمز واحد، نقرة قصيرة. إذن الرسالة كلها مكتوبة في الفواصل. فبأي لغة تتخاطب العصبونات؟

\section*{كم يمكن أن تقول النقرات}

لنبدأ كما يبدأ مهندس الاتصالات: كم من المعلومات يتسع له قناة كهذه أصلًا؟ إذا ميّز المستقبِل زمن وصول النقرة بدقة جزء من ألف من الثانية، أمكن للنقرة الواحدة أن تحمل نحو ثماني بتات. أما إذا اكتفى المستقبِل بعدّ النقرات التي وصلت في الثانية، فسيستخرج من التدفق نفسه أقل بعشرات المرات. سعة القناة كلها تقريبًا كامنة في \emph{الزمن}.

\pencilfig{4}{\pencilart{4}}{في الأعلى شيفرة مورس، وفي الأسفل عصبون: رمز واحد هو النقرة؛ والرسالة كلها مكتوبة في الفواصل.}

لكن القناة ضاجّة. والعصبون نفسه، على غير ما يُتوقع، دقيق: إذا غذّيناه مرات كثيرة بالتيار المتغير نفسه، نقر في اللحظات نفسها بدقة جزء من ألف من الثانية. الضجيج في الوصلات: أكثر من نصف النقرات التي تبلغ طرف السلك يضيع ببساطة، لأن المادة لا تُقذف في كل مرة. وفي القشرة الحية يبدو تدفق النقرات عشوائيًا تقريبًا. أهذا ضجيج، أم رسالة لا نعرف بعد كيف نقرؤها؟ الرسالة المضغوطة جيدًا لا يميزها الغريب عن العشوائية، فالسؤال مفتوح.

\section*{بطيء لكنه موثوق}

أبسط ترميز هو عدّ النقرات: كلما كثرت، كبر العدد. هذا الترميز لا يخشى الفقد ولا الارتعاش. لكنه بطيء. لمعرفة التردد بدقة عشرة في المئة نحتاج إلى مئة نقرة، أي إلى ثوانٍ. وفي المقابل تتعرف على حيوان في صورة عُرضت لك لحظة في نحو عُشر ونصف من الثانية، وتمر الإشارة في أثناء ذلك بنحو عشر مراحل معالجة. في كل مرحلة لا يلحق كل عصبون أن ينقر إلا مرة على الأكثر.

المخرج أن نأخذ عصبونات كثيرة معًا: لا مستمع واحد يعد النقرات دقيقة، بل مئة مستمع لحظة واحدة. لكن هنا أيضًا فخ: ضجيج العصبونات المتجاورة متشابه قليلًا، كالإشاعات في حشد. المتوسط على بضع عشرات من الجيران يساعد، ثم يكاد يتوقف عن المساعدة.

\pencilfig{122}{%
% error of the group-average rate relative to one neuron, sqrt(c + (1-c)/N): c = 0.1 (neighbours' noise
% is slightly alike; full edition, chapter 4) and c = 0 (independent noise). x = 0.8 log2 N, y = 2.2 * error
\draw[pen] (0,0) -- (5.9,0); \draw[pen] (0,0) -- (0,2.45);
\foreach \x/\n in {0/1, 2.66/10, 5.31/100}{\draw[pcGray, line width=0.5pt] (\x,0) -- (\x,-0.08); \node[lbl, anchor=north] at (\x,-0.08) {\n};}
\node[lbl, anchor=north] at (2.95,-0.38) {عدد العصبونات في المتوسط};
\node[lbl, anchor=south, rotate=90] at (-0.1,1.2) {الخطأ};
\draw[pcGray, densely dashed, line width=0.5pt] (0,0.70) -- (5.6,0.70);
\draw[penlite=pcBlue] plot[smooth] coordinates {(0.00,2.20) (0.47,1.80) (0.80,1.56) (1.27,1.27) (1.60,1.10) (2.07,0.90) (2.40,0.78) (2.87,0.64) (3.20,0.55) (3.67,0.45) (4.00,0.39) (4.47,0.32) (4.80,0.28) (5.27,0.22) (5.60,0.19)};
\draw[penlite=pcRed] plot[smooth] coordinates {(0.00,2.20) (0.47,1.84) (0.80,1.63) (1.27,1.39) (1.60,1.25) (2.07,1.10) (2.40,1.01) (2.87,0.92) (3.20,0.87) (3.67,0.82) (4.00,0.79) (4.47,0.76) (4.80,0.74) (5.27,0.73) (5.60,0.72)};
\node[lbl, text=pcRed, anchor=west] at (5.7,0.74) {ضجيج متشابه (كالقشرة)};
\node[lbl, text=pcBlue, anchor=west] at (5.7,0.19) {ضجيج مستقل};
\node[lbl, anchor=west] at (0.9,2.15) {عصبون واحد};
\node[lbl, anchor=north west] at (0.05,0.66) {أقل بثلاث مرات};
}{خطأ التردد المأخوذ متوسطًا على مجموعة من العصبونات. لو كان الضجيج مستقلًا، لواصل الهبوط. لكن تشابه ضجيج الجيران يوقفه عند الثلث تقريبًا، ويُبلغ هذا الحد بعد بضع عشرات من العصبونات فقط.}

\section*{سريع: من الأول}

المخرج الثاني أن نكتب العدد في الزمن. الإشارة القوية تجعل العصبون ينقر مبكرًا، والضعيفة متأخرًا. فالنقرة الواحدة تحمل عددًا. لكن من أين نبدأ عدّ الزمن إذا لم تكن هناك ساعة؟ يمكن أن نقارن العصبونات بعضها ببعض: من نقر أولًا، ومن ثانيًا. ترتيب وصول نقرات عشرة عصبونات يحمل أكثر من عشرين بتًا. ثم إن الرشقة الشبكية المشتركة تصلح علامة للبداية، كالفاصل الطويل بين الكلمات في شيفرة مورس.

\pencilfig{121}{%
% latency code: one click per neuron; the stronger the input, the earlier the click
\foreach \y/\t/\l/\k in {1.5/1.1/{مدخل قوي}/1, 0.95/2.4/{متوسط}/2, 0.4/4.2/{ضعيف}/3}{
  \draw[pen] (0.4,\y) -- (5.1,\y);
  \draw[pcBlue, line width=0.9pt] (\t,\y) -- (\t,\y+0.42);
  \node[lbl, anchor=east] at (0.3,\y+0.1) {\l};
  \draw[dotpen=pcAmber, wash=pcAmber] (5.6,\y+0.16) circle (0.17);
  \node[font=\scriptsize, text=pcAmber!60!black] at (5.6,\y+0.16) {\k};}
\draw[pcGray, densely dashed, line width=0.5pt] (0.55,0.25) -- (0.55,2.1);
\node[lbl, anchor=south] at (0.55,2.1) {المنبه};
\node[lbl, text=pcAmber!70!black, anchor=south] at (5.6,2.1) {الترتيب};
\draw[arr] (0.7,0.05) -- (4.9,0.05); \node[lbl, anchor=north] at (2.8,0.02) {الزمن};
}{مدخل قوي يعني نقرة مبكرة، وضعيف يعني نقرة متأخرة. ترتيب النقرات يخبر أي المداخل أقوى، حتى لو كانت لحظة البداية مجهولة.}

\section*{الترميز يختاره المستمع}

السلسلة نفسها من النقرات تحمل التردد والأزمنة والترتيب معًا. وأي منها سيُقرأ يقرره المستقبِل. العصبون ذو ”الثقب الواسع“ في دلوه يجمع الماء ببطء ولا يسمع إلا التردد. والعصبون السريع النسيان لا يسمع إلا التزامن. أما المكابح، كما رأينا، فتستطيع نقل العصبون من نمط إلى آخر.

حتى الوصلات ترشّح الإشارة. بعضها يتعب سريعًا، فلا ينقل التردد بل \emph{تغيّره}. وبعضها، على العكس، يمرر على نحو أفضل الرشقات المؤلفة من عدة نقرات متتالية: فتظهر لدى العصبون ”شرطة“. وفي الأثناء تضع العصبونات الكابحة علامات الترقيم: تقطّع التدفق إلى نوافذ وتخفضه إذا علا أكثر مما ينبغي.

\section*{لغة مقتصدة}

وأخيرًا: هذه اللغة غالية. كل نقرة تكلف طاقة، والدماغ كله يستهلك نحو عشرين واطًا. إذا قسمنا هذه الواطات على كل العصبونات، حصلنا في المتوسط على نحو نقرة واحدة في الثانية لكل عصبون. معظم العصبونات صامتة معظم الوقت. على ترميز الدماغ أن يكون بخيلًا.

في شيفرة مورس أكثر حروف الإنجليزية تواترًا، \foreignlanguage{english}{E}، نقطة واحدة: المتكرر يُرمَّز باختصار. والأمر مشابه عند العصبونات: المتوقع يكلف نقرات قليلة. المنبه الثابت يكف تدريجيًا عن إثارة الاستجابة، والمستشعرات تبلغ عن التغيرات، والدوبامين لا يبلغ إلا عن المفاجآت. الدماغ ينفق نقراته على الأخبار.

\fullversion{4}
